\section{Related work}%
\label{sec:related}

\paragraph{3D Reconstruction.}

There is a long and rich history of research on 3D reconstruction, beginning with seminal works that established the theory of multi-view geometry~\cite{faugeras90motion, oliensis00a-critique, hartley04multiple, ozyesil17a-survey}.
Follow-up work led to major practical advances, including robust SfM systems such as COLMAP~\cite{schonberger16structure-from-motion} and other pipelines~\cite{snavely06photo, agarwal11building, frahm10building, liu24robust, dellaert12factor}.
In this paper, we focus on \textit{feed-forward reconstruction models}, \ie, neural networks that infer scene geometry and camera poses directly from one or more images.
While recent SfM pipelines increasingly include learnable components such as keypoint detectors~\cite{yi16lift:,daniel18superpoint:,dusmanu19d2-net:,tyszkiewicz20disk:} and feature matchers~\cite{sarlin20superglue:,chen21learning,shi22clustergnn:,lindenberger23lightglue:}, our work is most closely related to end-to-end differentiable SfM frameworks that learn geometry estimation directly~\cite{zhou17unsupervised, ummenhofer17demon:, tang19ba-net:, wei20deepsfm:, wang21deep, teed20deepv2d:, teed21droid-slam:, brachmann24scene, wang23posediffusion:, wang24vggsfm:}.
While these works demonstrate that end-to-end learning is possible in SfM, they still \emph{combine} elements of classical SfM pipelines.
\duster~\cite{wang24dust3r:} and its follow-up \master~\cite{duisterhof25mast3r-sfm:} estimated both scene geometry and camera parameters (extrinsics and intrinsics) directly from images.
However, similar to stereo approaches~\cite{niemeyer20differentiable, fu22geo-neus:, wei21nerfingmvs:, yariv20multiview, yao18mvsnet:, gu20cascade, ma22multiview, peng22rethinking, zhang23geomvsnet:}, the neural networks in \duster and \master operate only on image pairs and still require post-optimization to process additional views.
A key improvement came from methods that process multiple images jointly, removing the need for optimization across views, including~\cite{wang24spann3r:, wang25continuous, tang25mv-dust3r:, wu25point3r:, zhang25flare:, yang25fast3r:, elflein25light3r-sfm:}.
Among these, VGGT~\cite{wang25vggt} is a representative approach that first surpassed post-optimization methods (\eg, using bundle adjustment) while relying solely on feed-forward inference, prompting many follow-up works~\cite{wang25p3, keetha26mapanything:, xiao25spatialtrackerv2:, chen25ttt3r:, deng25vggt-long:, lan26stream3r:, maggio25vggt-slam:, li25wint3r:, huang25longsplat:, szymanowicz26lagernvs, yuan25test3r:, deng26sail-recon:, jiang25anysplat:, cabon25must3r:, liu25plana3r:, li25rig3r:, zang26robust, chen26hd-vggt:, liu26streamcachevggt:, zhuo26streaming, jung26more:, mahdi25evict3r:, wang26lidar-vggt:, qi25fupad:, mazur264d-primitive-mache:}.
Several works improve the computational scalability of VGGT-style models through token merging, sparse attention, or descriptor-based aggregation, making many-view feed-forward reconstruction more efficient~\cite{shen26fastvggt:, wang25faster, mahdi25evict3r:}.
Other works address the quadratic cost of global self-attention by using stateful, bounded scene representations with a linear update cost.
Examples include methods that use test-time-training layers~\cite{chen26ttt3r:, jin26zipmap:, zhang26loger:}.
Others~\cite{zhuo26streaming, lan26stream3r:, chen26geometric, taher25kv-tracker:} process frames causally while maintaining a persistent geometric state by caching tokens, key-value pairs, local windows, or maps.
Some have used VGGT to reconstruct scenes locally, fusing results by means of alignment, localization, or SLAM algorithms~\cite{deng25vggt-long:, maggio25vggt-slam:, deng26sail-recon:, wang25amb3r:}.

Other works have extended VGGT-like models to take additional geometric information as input (\eg, cameras~\cite{keetha26mapanything:}), add sensors (\eg, LiDAR~\cite{wang26lidar-vggt:}), or model camera rigs~\cite{li25rig3r:}.
Other extensions predict normals~\cite{fang26dens3r:} or use mixture-of-experts designs~\cite{gao25more:, wang26moe3d:}.
PI3~\cite{wang25p3} removes the reliance on a fixed reference view, while Depth Anything 3 (DA3)~\cite{lin25depth} adopts vanilla DINO as its backbone, with both supporting dynamic scenes.
Some studies have investigated what these models learn, probing for correspondences, epipolar structure, and connections to 3D shape perception in humans~\cite{bratulic25on-geometric, bonnen26human-level}.
SelfEvo~\cite{huang26self-improving} explored a self-distillation scheme using spatiotemporal context asymmetry. 
Beyond reconstruction itself, the features learned by feed-forward reconstruction models have also been applied as geometry-aware representations for other tasks.
This includes video generation and novel view synthesis~\cite{jiang25anysplat:, huang26gen3r:, jang26repurposing, liu26camgeo:, szymanowicz26lagernvs}, vision-language models~\cite{zeng25janusvln:, qian26xembodied:, zhang26spatialstack:, zhang26make, zheng25learning, wang26let-geometry, yang26cambrian-s:, zhao25spacemind:, wu26mvggt:, lee253d-aware}, vision-language-action models~\cite{li25spatial, abouzeid25geoaware-vla:, vuong25improving, rao26augvla-3d:}, and perception tasks such as detection, segmentation, matching, occupancy prediction, and place recognition~\cite{cao26vggt-det:, gao26vggt-segmentor:, yang25dense, zhou26generalizing, xu26vggt-mpr:, deng25unipr-3d:}.

Monocular dynamic 3D reconstruction, or 4D reconstruction, aims to recover scene geometry that changes over time.
This line of research also has a long history, with early work by Bregler et al.~\cite{bregler00recovering} and Torresani et al.~\cite{torresani08nonrigid}.
Among recent contributions, MegaSaM~\cite{li25megasam:} has been particularly influential, combining feed-forward depth prediction with optimization-based non-rigid reconstruction.
ViPE~\cite{huang25vipe:} further builds on this direction.
Several works have explored feed-forward 4D reconstruction with reduced optimization.
MonST3R~\cite{zhang24monst3r:} and D$^2$USt3R~\cite{han25d2ust3r:} extend \duster to handle dynamic 3D content.
Align3R~\cite{lu25align3r:} builds on \duster to infer cameras and align monocular depth predictions over time, though it still relies on optimization beyond two views.
CUT3R~\cite{wang25continuous} and Point3R~\cite{wu25point3r:} support incremental reconstruction alongside dynamic scenes.
Geo4D~\cite{jiang25geo4d} fine-tunes a video generator to recover 4D geometry.
PI3~\cite{wang25p3} and DA3~\cite{lin25depth} adopt VGGT-style models and train them with dynamic data, while PAGE-4D~\cite{zhou25page-4d:} adapts VGGT through a module that separates static and moving regions.
Human3R focuses on human-scene reconstruction~\cite{chen25human3r:}.
SpatialTracker~\cite{xiao25spatialtrackerv2:}, St4RTrack~\cite{feng25st4rtrack:}, DPM~\cite{sucar25dynamic}, V-DPM~\cite{sucar26v-dpm}, Uni4D~\cite{yao25uni4d:}, Any4D~\cite{karhade2025any4d}, and related methods unify dynamic reconstruction and point tracking.
D4RT~\cite{zhang25efficiently} introduces a unified feed-forward transformer that reconstructs dynamic scenes by querying point-level 4D scene information.


\paragraph{Registers in Vision Transformers (ViTs).}

Recent works~\cite{darcet24vision, marouani26revisiting} have found that a small number of tokens in ViTs encode not only local patch information but also global information.
These outliers have high norms and disrupt the spatial coherence of the patch features.
This, in turn, has motivated the use of register tokens, which are separate from image tokens and help preserve coherence among patch representations. 
Further studies have investigated the causes of this phenomenon and the interactions between register and image tokens~\cite{jiang25vision, shi26vision, lappe25register, marouani26revisiting}.
Here, we add register tokens to each input image frame and use them to aggregate and exchange global information across images.
In this way, our registers carry information about the sequence as a whole.
Rather than discarding them as is often done in prior work, we show that they are useful in downstream applications.
